%% ============================================================
%%  Anexo C – Árbol de Potencia Detallado
%% ============================================================
\chapter{Árbol de Potencia Detallado}
\label{chap:anexo_potencia}

\begin{table}[H]
  \centering
  \caption{Rieles de alimentación del sistema ALMA -- detalle por riel.}
  \label{tab:potencia_detalle}
  \begin{tabularx}{\textwidth}{L{2cm} C{1.5cm} C{1.5cm} C{1.8cm} C{1.8cm} L{2.5cm} L{2.5cm}}
    \toprule
    \rowcolor{udea-green}
    \textcolor{white}{\textbf{Riel}} &
    \textcolor{white}{\textbf{V nom.}} &
    \textcolor{white}{\textbf{Tol.}} &
    \textcolor{white}{\textbf{I típica}} &
    \textcolor{white}{\textbf{I pico}} &
    \textcolor{white}{\textbf{Generado por}} &
    \textcolor{white}{\textbf{Consumido por}} \\
    \midrule
    VIN      & 12\,V   & $\pm$5\,\%  & 1.6\,A    & 2.5\,A    & USB C & PCB-PWR (todos los Buck) \\
    \rowcolor{row-alt}
    VDD\_NPU & 1.1\,V  & $\pm$2\,\%  & 1.0\,A    & 2.0\,A    & SY88003 (Silergy) & Núcleo NPU A311D2 \\
    VDD\_CPU & Variable & $\pm$2\,\% & 2.0\,A    & 4.0\,A    & MPU8756 (MPS)     & CPU-A, CPU-B A311D2 \\
    \rowcolor{row-alt}
    VDD\_IO  & 1.8\,V  & $\pm$2\,\%  & 0.5\,A    & 1.0\,A    & TL76701 (TI)      & E/S SoC, buses lógicos \\
    VCC5     & 5.0\,V  & $\pm$3\,\%  & 0.8\,A    & 2.5\,A    & NB679 (MPS)       & Amp. Clase~D, Boost backlight \\
    \rowcolor{row-alt}
    VCC3V3   & 3.3\,V  & $\pm$3\,\%  & 0.6\,A    & 1.0\,A    & NB680 (MPS)       & WiFi/BT, eMMC, GPIO \\
    VLED     & 5.0\,V  & ---         & 0.3\,A    & 1.1\,A    & NB679 (compartido) & LEDs RGB \\
    \bottomrule
  \end{tabularx}
\end{table}

\textit{Nota: los valores de corriente típica y pico son estimaciones basadas en
los datasheets de los componentes seleccionados y en el presupuesto energético del
Entregable~2 (19\,W totales). Los anchos de pista se calcularon con IPC-2152
\cite{ipc2152} -- ver Anexo~\ref{chap:anexo_si}, Sección~B.3.}

\section*{C.1 Presupuesto Energético Total}

Siguiendo la estimación del Entregable~2, el consumo máximo del sistema en
condiciones de carga pico es de aproximadamente 19\,W. Con la entrada de 12\,V:

\[
  I_{total} \approx \frac{19\,\text{W}}{12\,\text{V}} \approx 1.6\,\text{A}
\]

El adaptador externo seleccionado debe tener una corriente nominal mínima de 2\,A
(24\,W) para proveer un margen de seguridad del $\approx$20\,\% sobre la carga
máxima, tal como se definió en el Entregable~2. Se recomienda un adaptador de 2.5\,A
para absorber los picos transitorios del amplificador Clase~D.

\section*{C.2 Notas sobre la Red de Feedback del SY88003}

El regulador SY88003 se utiliza en múltiples instancias del esquemático de la
PCB-PWR, cada una configurada para un voltaje de salida diferente mediante la red
de resistencias de feedback (RFB). La relación entre los valores de la red y el
voltaje de salida es:

\[
  V_{out} = V_{ref} \times \left(1 + \frac{R_{top}}{R_{bot}}\right)
\]

donde $V_{ref}$ es el voltaje de referencia interno del SY88003 (verificar en
datasheet). Los valores exactos de $R_{top}$ y $R_{bot}$ para cada instancia
(NPU~1.1\,V, GPU, CPU, etc.) están tomados del esquemático de referencia
de Khadas VIM4 y deben preservarse sin modificación para garantizar la
estabilidad del SoC.
